\chapter{Twierdzenia o rzędzie}
\label{ma2:ch:rank}

W tym rozdziale badamy lokalną postać odwzorowań,
których pochodna ma stały rząd. W skończonym wymiarze
wykorzystamy macierze Jacobiego, a następnie opiszemy
warunki pozwalające przeprowadzić analogiczną konstrukcję
w przestrzeniach Banacha.

\section{Rząd i jego lokalna stałość}

\begin{definicja}
\label{ma2:def:rank}
Niech $U\subseteq\R^m$ będzie otwarty i
$f:U\to\R^n$ różniczkowalne. Rzędem $f$ w punkcie
$p\in U$ nazywamy $\operatorname{rz}Df(p)$, czyli
wymiar obrazu operatora $Df(p)$. Jest to rząd
macierzy Jacobiego w bazach standardowych.
Jeżeli $Df(p)$ jest różnowartościowy, mówimy
o immersji w punkcie $p$; jeżeli jest surjektywny,
mówimy o submersji w punkcie $p$.
\end{definicja}

\begin{stwierdzenie}
\label{ma2:prop:rank-lower}
Jeżeli $f\in C^1(U,\R^n)$ i
$\operatorname{rz}Df(p)=r$, to w pewnym otoczeniu $p$
zachodzi $\operatorname{rz}Df(x)\ge r$.
W szczególności immersja lub submersja w jednym
punkcie pozostaje nią w pewnym otoczeniu tego punktu.
\end{stwierdzenie}
\begin{proof}
Dla $r=0$ nierówność jest automatyczna. Dla $r>0$
istnieje niezerowy minor stopnia $r$ macierzy $Df(p)$.
Wyznacznik tego minora, obliczany dla $Df(x)$, jest
ciągły jako wielomian w ciągłych współczynnikach.
Jego wartość pozostaje niezerowa w otoczeniu $p$:
wystarczy, by różnica od wartości w $p$ miała moduł
mniejszy od połowy modułu tej wartości.
Rząd macierzy jest więc co najmniej $r$.
W przypadku immersji $r=m$, a w przypadku submersji
$r=n$; rząd nie może przekraczać $\min\{m,n\}$,
więc w obu przypadkach pozostaje maksymalny.
\end{proof}

\section{Twierdzenie o rzędzie stałym}

\begin{twierdzenie}
\label{ma2:thm:constant-rank}
Niech $f:U\subseteq\R^m\to\R^n$ będzie klasy $C^s$,
gdzie $s\ge1$ lub $s=\infty$. Załóżmy, że
$\operatorname{rz}Df(x)=r$ w pewnym otwartym otoczeniu
$p\in U$. Istnieją lokalne dyfeomorfizmy $\alpha,\beta$
klasy $C^s$ w otoczeniach odpowiednio $p$ i $f(p)$,
przeprowadzające te punkty na zera, takie że
\[
 (\beta\circ f\circ\alpha^{-1})(u,v)=(u,0),
 \qquad (u,v)\in A\times B,
\]
gdzie $A\subseteq\R^r$, $B\subseteq\R^{m-r}$
są otwartymi otoczeniami zera. Przeciwdziedzinę
identyfikujemy z $\R^r\times\R^{n-r}$.
\end{twierdzenie}

Dla klasy $C^\infty$ jest to szczególny przypadek
\atlasref{DG-01::tw:rzad-staly}{DG-01: twierdzenia o rzędzie stałym}.
Konstrukcję dla skończonego $s$ przedstawimy za pomocą
poniższego wyniku dla przestrzeni Banacha. Warunek
zerowania pewnego operatora w tej konstrukcji zastąpi
argument z minorami użyty w dowodzie w DG-01.

\section{Postać blokowa w przestrzeniach Banacha}

Niech $E=E_1\oplus E_0$ i $F=F_1\oplus F_0$
będą topologicznymi sumami prostymi domkniętych
podprzestrzeni. Oznacza to, że odwzorowania sumowania
z produktów na $E,F$ są izomorfizmami liniowymi
i topologicznymi. Używamy współrzędnych produktowych
oraz norm maksimum, równoważnych normom wyjściowym.

\begin{twierdzenie}[Postać blokowa]
\label{ma2:thm:rank-banach-block}
Niech $f=(f_1,f_0):U\subseteq E_1\times E_0
\to F_1\times F_0$ będzie klasy $C^s$, $s\ge1$
lub $s=\infty$, gdzie $U$ jest otwarty i zawiera $0$,
$f(0)=0$. Zapiszmy
\[
 Df(x)=
 \begin{pmatrix} A(x)&B(x)\\ C(x)&D(x)\end{pmatrix}.
\]
Załóżmy, że $A(0):E_1\to F_1$ jest izomorfizmem
liniowym i topologicznym i w pewnym otoczeniu zera
\[
 D(x)-C(x)A(x)^{-1}B(x)=0.
\]
W tym warunku rozpatrujemy otoczenie, na którym
$A(x)$ jest odwracalny. Wtedy istnieją lokalne
dyfeomorfizmy klasy $C^s$, sprowadzające $f$
do postaci $(u,v)\mapsto(u,0)$ na otoczeniu
produktowym w $F_1\times E_0$.
\end{twierdzenie}
\begin{proof}
Wszystkie podprzestrzenie są przestrzeniami Banacha,
jako domknięte podprzestrzenie przestrzeni Banacha.
Ciągłość $A(x)$ w normie operatorowej pozwala
zmniejszyć otoczenie zera tak, by
$\|A(0)^{-1}(A(x)-A(0))\|<1$.
Z twierdzenia \ref{ma2:thm:inverse-perturbation}
w algebrze $\mathcal L(E_1)$ wynika odwracalność
$A(0)^{-1}A(x)$, a więc również $A(x)$.
Jeśli $E_1=\{0\}$, to także $F_1=\{0\}$ i
odwracalność jest automatyczna.

Niech $\alpha(x_1,x_0)=(f_1(x_1,x_0),x_0)$.
Jej pochodna w zerze przeprowadza $(h,k)$ na
$(A(0)h+B(0)k,k)$. Odwrotny operator przeprowadza
$(u,v)$ na $(A(0)^{-1}(u-B(0)v),v)$.
Podstawienie w obu kierunkach daje identyczność,
a ograniczoność wynika z ograniczoności $A(0)^{-1}$
i $B(0)$ oraz normy maksimum. Z twierdzenia
\ref{ma2:thm:inverse-function} otrzymujemy lokalny
dyfeomorfizm $\alpha$ klasy $C^s$.
Ograniczamy obraz $\alpha$ do otoczenia produktowego
$P\times Q\subseteq F_1\times E_0$ złożonego z kul
otwartych o środkach w zerze, a jego dziedzinę
do przeciwobrazu tego produktu.

Niech $g=f\circ\alpha^{-1}$. Z konstrukcji
$g(u,v)=(u,H(u,v))$. Oznaczając $x=\alpha^{-1}(u,v)$,
z reguły łańcuchowej i wzoru na pochodną odwrotności
otrzymujemy
\[
 Dg(u,v)=
 \begin{pmatrix}A&B\\C&D\end{pmatrix}
 \begin{pmatrix}A^{-1}&-A^{-1}B\\0&I_{E_0}\end{pmatrix}
 =
 \begin{pmatrix}I_{F_1}&0\\CA^{-1}&D-CA^{-1}B\end{pmatrix},
\]
gdzie wszystkie bloki po prawej stronie są obliczone
w $x$. Z założenia $D_vH=0$.
Dla $u\in P$, $v\in Q$ odcinek $(u,tv)$,
$0\le t\le1$, leży w produkcie.
Funkcja $t\mapsto H(u,tv)$ ma zerową pochodną,
więc jest stała na mocy nierówności o wartości
średniej (twierdzenie \ref{ma2:thm:mean-value-inequality}).
Zatem $H(u,v)=H(u,0)=h(u)$, gdzie $h$ jest klasy $C^s$.

Na $P\times F_0$ definiujemy
$\beta(u,w)=(u,w-h(u))$. Odwzorowanie jest klasy
$C^s$ i ma odwrotność $(u,z)\mapsto(u,z+h(u))$
tej samej klasy. Obie kompozycje dają identyczność,
a dziedzina i obraz są otwarte. Ponieważ $h(0)=0$,
$\beta(0)=0$. Ostatecznie
$\beta\circ f\circ\alpha^{-1}(u,v)=(u,0)$.
\end{proof}

\begin{proof}[Dowód twierdzenia \ref{ma2:thm:constant-rank}]
Przez przesunięcia zakładamy $p=0$, $f(p)=0$.
Gdy $r>0$, wybieramy niezerowy minor stopnia $r$,
a przez permutacje współrzędnych umieszczamy go
w lewym górnym bloku $A(0)$ macierzy $Df(0)$.
Ciągłość zapewnia odwracalność $A(x)$ w otoczeniu zera.
Dla każdego takiego $x$ zachodzi
\[
 \begin{pmatrix}I&0\\-CA^{-1}&I\end{pmatrix}
 \begin{pmatrix}A&B\\C&D\end{pmatrix}
 \begin{pmatrix}I&-A^{-1}B\\0&I\end{pmatrix}
 =
 \begin{pmatrix}A&0\\0&D-CA^{-1}B\end{pmatrix}.
\]
Macierze skrajne są odwracalne (zmiana znaku
bloku pozadiagonalnego daje odwrotności).
Mnożenie przez nie zachowuje rząd. Rząd prawej
strony wynosi $r+\operatorname{rz}(D-CA^{-1}B)$,
gdyż obrazy bloków leżą w różnych składowych
sumy prostej. Ponieważ rząd $Df(x)$ jest równy $r$,
drugi blok musi być zerowy. Stosujemy twierdzenie
\ref{ma2:thm:rank-banach-block}.
Dla $r=0$ pochodna jest zerowa w otoczeniu zera;
na mniejszej kuli $f$ jest stałe z nierówności
o wartości średniej. Przesunięcia dają wtedy
żądaną postać zerową. Jeśli $r=m$ lub $r=n$,
brakujący blok ma zerową dziedzinę lub przeciwdziedzinę
i powyższy argument pozostaje ważny.
\end{proof}

W nieskończonym wymiarze sama równość wymiarów obrazów
nie gwarantuje zerowania bloku resztowego.
Na przykład na $\ell^2\times\R$ odwzorowanie
$(u,t)\mapsto(u,t^2)$ ma nieskończenie wymiarowy
obraz pochodnej w każdym punkcie, lecz jądro pochodnej
ma wymiar $1$ dla $t=0$ i $0$ dla $t\ne0$.
Nie może więc mieć lokalnej postaci o stałym jądrze
w otoczeniu $(0,0)$: reguła łańcuchowa dla
dyfeomorfizmów zachowuje wymiar jądra pochodnej.

\section{Immersje, submersje i poziomice}

\begin{wniosek}
\label{ma2:cor:rank-maximal}
Immersja w $p$ ma lokalną postać $u\mapsto(u,0)$,
a submersja w $p$ ma lokalną postać $(u,v)\mapsto u$,
w odpowiednich współrzędnych klasy $C^s$.
\end{wniosek}
\begin{proof}
Stwierdzenie \ref{ma2:prop:rank-lower} daje stałość
maksymalnego rzędu w otoczeniu $p$.
Stosujemy twierdzenie \ref{ma2:thm:constant-rank}
z $r=m$ albo $r=n$. W pierwszym przypadku
$\R^{m-r}=\{0\}$, w drugim $\R^{n-r}=\{0\}$.
\end{proof}

Są to euklidesowe odpowiedniki
\atlasref{DG-01::wn:immersja-lokalnie}{DG-01: lokalnej postaci immersji}
oraz
\atlasref{DG-01::wn:submersja-lokalnie}{DG-01: lokalnej postaci submersji}.

\begin{definicja}
\label{ma2:def:regular-value}
Wartość $c\in\R^n$ nazywamy regularną dla
$f:U\to\R^n$ klasy $C^1$, jeżeli $Df(p)$ jest
surjektywny dla każdego $p\in f^{-1}(\{c\})$.
Wartość spoza obrazu jest regularna w sensie
pustego warunku.
\end{definicja}

\begin{wniosek}
\label{ma2:cor:regular-fiber}
Niech $c$ będzie wartością regularną odwzorowania
$f$ klasy $C^s$. W każdym punkcie niepustej poziomicy
$f^{-1}(\{c\})$ istnieją współrzędne lokalne,
w których poziomica ma postać $\{0\}\times B$,
gdzie $B$ jest otwarty w $\R^{m-n}$.
\end{wniosek}
\begin{proof}
Dla $p$ na poziomicy mamy $m\ge n$ i submersję
w $p$. Z wniosku \ref{ma2:cor:rank-maximal} otrzymujemy
$\beta(f(\alpha^{-1}(u,v)))=u$, gdzie $\beta(c)=0$.
Ponieważ $\beta$ jest różnowartościowa, równanie
$f(\alpha^{-1}(u,v))=c$ jest równoważne $u=0$.
Ograniczenie dziedziny współrzędnych do produktu
otwartych kul daje tezę.
\end{proof}

W następnym rozdziale wykorzystamy tę postać do
opisu podrozmaitości i ich przestrzeni stycznych.

\section{Przykład geometryczny}

\begin{przyklad}[Submersja i jej poziomice]
\label{ma2:ex:rank-paraboloids}
Niech $f:\R^3\to\R$ będzie dane wzorem
$f(x,y,z)=z-x^2-y^2$. Mamy
$Df(x,y,z)(h,k,l)=l-2xh-2yk$.
Operator ten jest surjektywny, ponieważ
$Df(x,y,z)(0,0,l)=l$ dla dowolnego $l\in\R$.
Każda wartość $c$ jest regularna, a jej poziomica
to paraboloida $z=x^2+y^2+c$. Globalna zmiana
współrzędnych $(x,y,z)\mapsto(z-x^2-y^2,x,y)$
ma odwrotność $(u,v,w)\mapsto(v,w,u+v^2+w^2)$:
obie kompozycje są identycznościami, a oba
odwzorowania są wielomianowe. W tych współrzędnych
$f$ jest rzutem na pierwszą składową.
\end{przyklad}

\paragraph{Rysunek 3D: rodzina poziomic submersji.}
Przedstawiamy trzy rozłączne fragmenty poziomic
dla $c=-1,0,1$ nad kwadratem $[-1,1]^2$.
\begin{lstlisting}[language=Python]
import numpy as np
import plotly.graph_objects as go

x = np.linspace(-1, 1, 61)
X, Y = np.meshgrid(x, x)
fig = go.Figure()
for c, palette in [(-1, "Blues"), (0, "Greens"),
                   (1, "Oranges")]:
    fig.add_trace(go.Surface(x=X, y=Y, z=X**2+Y**2+c,
        colorscale=palette, showscale=False,
        opacity=0.8, name=f"c={c}", showlegend=True))
fig.update_layout(scene=dict(xaxis_title="x",
    yaxis_title="y", zaxis_title="z"))
fig.write_html("am2_20_rank.html", include_plotlyjs=True)
fig.show()
\end{lstlisting}

\begin{przyklad}[Zmiana rzędu]
\label{ma2:ex:rank-jump}
Dla $f(x,y)=(x,y^2)$ macierz Jacobiego jest
diagonalna o wpisach $1,2y$. Rząd wynosi $1$
dla $y=0$ oraz $2$ dla $y\ne0$.
W żadnym otoczeniu punktu osi $y=0$ nie ma
postaci normalnej o stałym rzędzie: pochodne
dyfeomorfizmów są izomorfizmami, więc reguła
łańcuchowa zachowuje rząd w każdym punkcie.
\end{przyklad}
